\documentclass[10pt]{article}
\usepackage[a4paper,margin=1in]{geometry}
\usepackage{amsmath,amssymb,amsthm}
\usepackage[hidelinks]{hyperref}
\newtheorem{proposition}{Proposition}
\newcommand{\R}{\mathbb R}
\title{An anisotropic Gaussian counterexample to the\\observable-gradient coupling conjecture}
\author{MathDB problem 339682}
\date{12 September 2026}
\begin{document}
\maketitle

\paragraph{Task 1.} Resolve the claimed sign of the variance change for
couplings built from the observable gradient. This is Conjecture 62,
equation (157), in the final publication of N\"usken and Pavliotis~\cite{NP}.
The arXiv version numbers it Conjecture 60.

For the overdamped generator \(L=-\nabla V\cdot\nabla+\Delta\), the source
requires centered \(f,\phi\in L^2(\pi)\) satisfying \(-L\phi=f\), where
\(\pi(dx)=Z^{-1}e^{-V(x)}dx\). Its conjectured inequality is
\begin{equation}\label{conj}
 \int_{\R^d\times\R^d}
 \nabla\phi(y)^T\alpha_f(x,y)\nabla\phi(x)
 e^{-V(x)-V(y)}\,dx\,dy\leq0.
\end{equation}
Writing \(\widehat{\nabla f}=\nabla f/|\nabla f|\), equation (156) of the source
sets
\[
 \alpha_f(x,y)=I-2\,\frac{w(x,y)w(x,y)^T}{|w(x,y)|^2},
 \qquad w(x,y)=\widehat{\nabla f}(x)+\widehat{\nabla f}(y),
\]
when both gradients and their normalized sum are nonzero; otherwise it sets
\(\alpha_f=I\). Admissibility in equation (151) is
\(\alpha_f^T\alpha_f\leq I\).

\begin{proposition}
Inequality \eqref{conj} is false for the smooth Gaussian potential
\[
 V(x)=\tfrac12(x_1^2+9x_2^2),\qquad f(x)=x_1+3x_2,\qquad d=2.
\]
Its left-hand side equals \(56\pi^2/405>0\).
\end{proposition}

\begin{proof}
Here \(Z=2\pi/3\) and \(\pi\) is the centered Gaussian with covariance
\(\operatorname{diag}(1,1/9)\). Put
\[
 \phi(x)=x_1+\tfrac13x_2,\qquad
 v=\nabla f=\begin{pmatrix}1\\3\end{pmatrix},\qquad
 a=\nabla\phi=\begin{pmatrix}1\\1/3\end{pmatrix}.
\]
Both linear functions are centered and square-integrable under \(\pi\).
Since their second derivatives vanish,
\[
 -L\phi
 =x_1\partial_1\phi+9x_2\partial_2\phi-\Delta\phi
 =x_1+3x_2=f.
\]
The observable gradient is the same nonzero vector at every point, so
\(w=2v/\sqrt{10}\neq0\). Thus the nondegenerate branch of the source's
definition applies everywhere and gives
\[
 \alpha_f=I-\frac{2vv^T}{v^Tv}
 =\frac15\begin{pmatrix}4&-3\\-3&-4\end{pmatrix}.
\]
This constant matrix is measurable and satisfies \(\alpha_f^T\alpha_f=I\).
The integrand apart from its Gaussian weight is therefore the constant
\[
 a^T\alpha_fa
 =|a|^2-\frac{2(a^Tv)^2}{|v|^2}
 =\frac{10}{9}-\frac{8}{10}
 =\frac{14}{45}>0.
\]
Multiplication by \(Z^2=4\pi^2/9\) gives \(56\pi^2/405>0\), as claimed.
If the integral is expressed using the probability measure
\(\pi\otimes\pi\), its value is \(14/45\); either normalization contradicts
the proposed sign.
\end{proof}

The example also defines genuine small perturbations of independent
sampling. For \(0<\varepsilon<1\), the constant diffusion matrix
\[
 Q_\varepsilon=
 \begin{pmatrix}I&\varepsilon\alpha_f\\
                 \varepsilon\alpha_f^T&I\end{pmatrix}
\]
has eigenvalues \(1-\varepsilon\) and \(1+\varepsilon\), each twice.
With drift matrix
\(D=\operatorname{diag}(1,9,1,9)\), the process
\(dU_t=-DU_t\,dt+\sqrt{2Q_\varepsilon}\,dW_t\)
has exactly the required two overdamped marginal processes.
Its covariance converges to
\(2\int_0^\infty e^{-Ds}Q_\varepsilon e^{-Ds}\,ds\), which is positive
definite, and its mean converges to zero. It is therefore the usual
nondegenerate ergodic Gaussian coupling. The counterexample does not
depend on using only a possibly degenerate full-strength endpoint.

The source's motivating integration-by-parts identity remains valid:
\(\int\nabla f\cdot\nabla\phi\,d\pi=2=\int f^2\,d\pi\).
That identity alone does not force the reflection built from
\(\nabla f\) to have a nonpositive quadratic form on \(\nabla\phi\).
Anisotropy of the potential separates these two constant directions.
The argument resolves the universal conjecture by a counterexample; it
does not decide its restrictions to one dimension or to isotropic targets.

\paragraph{Task 2.} Determine the sharp condition-number threshold for
nonpositive variance change for every linear observable under a Gaussian
target, and classify the extremal directions.

\paragraph{Task 3.} Decide whether the observable-gradient sign conjecture
holds in one dimension under smooth uniform convexity of the potential,
including observables whose derivatives change sign.

\paragraph{Task 4.} Construct practical preconditioners for the observable
gradient and prove quantitative conditions under which the resulting
reflection guarantees a variance improvement.

\paragraph{Task 5.} Quantify the largest variance increase, relative to
independent sampling, for finite coupling strengths in Gaussian targets
with prescribed covariance condition number.

\begin{thebibliography}{9}
\raggedright
\bibitem{NP}
N.~N\"usken and G.~A.~Pavliotis,
\emph{Constructing Sampling Schemes via Coupling: Markov Semigroups and
Optimal Transport}, SIAM/ASA Journal on Uncertainty Quantification
\textbf{7}(1) (2019), 324--382.
\href{https://doi.org/10.1137/18M119896X}{doi:10.1137/18M119896X}.
Section 5.2, pp.~364--365, equations (150)--(157).
\href{https://www.ma.imperial.ac.uk/~pavl/NuskenPavliotis2018.pdf}{Final author-hosted publisher PDF}.
\end{thebibliography}
\end{document}
